\subsection{theoryv3 Evidence Layer (Compressed, Auto-Generated)}
\label{sec:theoryv3-evidence}

The \texttt{theoryv3} analysis branch is a compact ``evidence layer'' for the discrete kernel and crosslink logic. Current inventory (auto-generated from the replication bundle; theoryv3 is a baseline branch pinned to v2.5 by design):
\[
\text{modules}=\TFPTTheoryVThreeModulesTotal,\qquad
\text{checks}=\TFPTTheoryVThreeChecksTotal\ (\mathrm{PASS/WARN/FAIL}=\TFPTTheoryVThreeChecksPassTotal/\TFPTTheoryVThreeChecksWarnTotal/\TFPTTheoryVThreeChecksFailTotal),\qquad
\text{plots expected/present}=\TFPTTheoryVThreePlotsExpectedTotal/\TFPTTheoryVThreePlotsPresentTotal.
\]

% Inlined (standalone build): formerly \input{../tfpt-suite/2.5/theoryv3/out/theoryv3_manifest_table.tex}
% -----------------------------------------------------------------------------
% AUTO-GENERATED by tfpt-suite/2.5/theoryv3/theoryv3_manifest.py — do not edit by hand.
\begin{table}[H]
\centering
\small
\begin{tabular}{@{}lllcc@{}}
\toprule
Module & Signal & Result & Checks & Plots OK\\
\midrule
\texttt{alpha\_backreaction\_sensitivity\_audit} & $\alpha$ backreaction sensitivity (k-grid) & PASS & 2/2 & yes\\
\texttt{axion\_dm\_audit} & Axion DM audit & PASS & 2/2 & yes\\
\texttt{baryon\_consistency\_audit} & Baryon consistency / H0 audit & PASS & 5/5 & yes\\
\texttt{cfe\_coefficient\_minimality\_audit} & cfe\_coefficient\_minimality\_audit & PASS & 3/3 & yes\\
\texttt{constant\_factory\_audit} & Constant factory audit (derived outputs) & PASS & 3/3 & yes\\
\texttt{dark\_energy\_exponential\_audit} & Dark energy exponential candidate & PASS & 2/2 & yes\\
\texttt{dark\_energy\_norm\_half\_origin\_audit} & Dark energy normalization n=1/2 scan & PASS & 2/2 & yes\\
\texttt{defect\_partition\_g5\_audit} & Two-defect partition ($\delta_2=(5/4)\delta_{\mathrm{top}}^2$) & PASS & 4/4 & yes\\
\texttt{flavor\_pattern\_audit} & Flavor pattern audit (texture invariants) & PASS & 3/3 & yes\\
\texttt{g5\_bridge\_defect\_to\_holonomy\_audit} & g5\_bridge\_defect\_to\_holonomy\_audit & PASS & 3/3 & yes\\
\texttt{g5\_crosslink\_audit} & Crosslink consistency (g=5 anchors) & PASS & 3/3 & yes\\
\texttt{g5\_origin\_audit} & Holonomy degeneracy g=5 & PASS & 2/2 & yes\\
\texttt{pmns\_tm1\_audit} & PMNS TM1 audit & PASS & 2/2 & yes\\
\texttt{seed\_invariants\_audit} & Seed invariants ($c_3, \varphi_0, \delta_{\mathrm{top}}, \beta$) & PASS & 7/7 & yes\\
\texttt{yukawa\_exponent\_index\_audit} & Yukawa exponent-index mapping & PASS & 3/3 & yes\\
\texttt{yukawa\_index\_mapping\_audit} & Yukawa index mapping & PASS & 3/3 & yes\\
\bottomrule
\end{tabular}
\caption{theoryv3 compressed evidence layer (auto-generated from the replication bundle).}
\end{table}

\begin{tcolorbox}[colback=gray!5, colframe=gray!60,
title={\textbf{Two ``lock-in'' results (negative controls; theoryv3)}}, fonttitle=\bfseries\small]
\textbf{$k=2$ is selected by a sensitivity minimum:} the \texttt{alpha\_backreaction\_sensitivity\_audit} shows the minimum $|\mathrm{ppm}|$ in the $k$ grid occurs at $k=2$, with $|\mathrm{ppm}|\simeq 0.037$.

\textbf{$g=5$ is selected by a defect-partition negative control:} the \texttt{defect\_partition\_g5\_audit} shows $g=5$ minimizes $|z|$ sharply versus $g\in\{4,6\}$ (gap $\sim 44.98$ in $|z|$).
\textbf{Evidence:} \texttt{defect\_partition\_g5\_audit}, \texttt{two\_defect\_cluster\_expansion\_audit}, \texttt{g5\_independence\_audit}, \texttt{g5\_trace\_completeness\_audit}.
\end{tcolorbox}

\begin{figure}[H]
\centering
\begin{tikzpicture}
\begin{groupplot}[
    group style={group size=2 by 1, horizontal sep=1.2cm},
    width=0.48\textwidth,
    height=0.42\textwidth,
    grid=both,
    legend style={font=\scriptsize},
]
\nextgroupplot[
    title={$k$-sensitivity (abs ppm)},
    xlabel={$k$},
    ylabel={$|{\rm ppm}|$},
    xmin=-0.1, xmax=3.1,
    ymin=0, ymax=4.2,
    xtick={0,1,1.5,2,2.5,3},
]
\addplot+[mark=*, mark size=2.2pt] coordinates {
    (0.0, 3.665371791638453)
    (1.0, 1.8074246412353618)
    (1.5, 0.883514212976573)
    (2.0, 0.03703710120260405)
    (2.5, 0.9542414840493093)
    (3.0, 1.8681110743677323)
};
\addplot[gray, dashed] coordinates {(2.0,0.0) (2.0,4.2)};

\nextgroupplot[
    title={$g$ negative control (z-score)},
    xlabel={$g$},
    ylabel={$z$},
    xmin=3.7, xmax=6.3,
    ymin=-60, ymax=60,
    xtick={4,5,6},
]
\addplot+[mark=*, mark size=2.2pt] coordinates {
    (4, -46.8455434804275)
    (5, 1.8646873709648437)
    (6, 50.5749164827739)
};
\addplot[gray, dashed] coordinates {(5,-60) (5,60)};
\end{groupplot}
\end{tikzpicture}
\caption{theoryv3 lock-in diagnostics: $k$-sensitivity of $\alpha^{-1}(0)$ (left) and the $g\in\{4,5,6\}$ negative control selecting $g=5$ (right).}
\end{figure}

\begin{tcolorbox}[colback=green!3, colframe=green!50!black,
title={\textbf{Global likelihood (v2.8)}}, fonttitle=\bfseries\small]
v2.8 introduces a dedicated \texttt{global\_likelihood\_v28} module that replaces the v2.7 dashboard $\chi^2$ with explicit covariance handling and dataset wiring.

\textbf{Score (10 enabled datasets):}
\[
\chi^2_{\mathrm{total}}=26.969\ (25~\mathrm{dof}),\qquad p\simeq 0.357,\qquad \mathrm{coverage\_ratio}=1.0.
\]
The core consistency test (excluding $\alpha$ metrology dominance):
\[
\chi^2_{\mathrm{core}}=2.437\ (6~\mathrm{dof}),\qquad p\simeq 0.875.
\]

\textbf{Disabled datasets (3, diagnostic only --- \emph{not} included in $\chi^2_{\mathrm{total}}$):}
\begin{itemize}[leftmargin=1.5em, topsep=2pt, itemsep=1pt]
    \item \texttt{H0\_planck2018} and \texttt{H0\_sh0es\_2022}: $H_0$ is derived as a diagnostic
          (Sec.~\ref{sec:known-gaps}; module \texttt{h0\_diagnostic\_closure\_v28}); additionally, the
          Hubble tension between these two references makes including either in the unconditional score misleading.
          $H_0$ is reported but not scored.
    \item \texttt{axion\_haloscope\_frequency\_target}: the axion frequency prediction is \emph{conditional}
          on the adopted axion scenario and is therefore scored separately in the conditional block.
\end{itemize}
This scoring policy is declared in the \texttt{global\_likelihood\_v28} dataset specification and is not adjusted post-hoc.

All 235 modules PASS in this snapshot; 0~WARN, 0~FAIL.

\textbf{Evidence:} \texttt{global\_likelihood\_v28} (PASS), \texttt{global\_consistency\_test} (PASS).
\end{tcolorbox}

\paragraph{$H_0$ consolidation policy.}
The suite computes $H_0$ through 12 internal pathways (spread ${\sim}2.4\%$). The
\textbf{authoritative value} reported in the PRL claim surface is
$H_0 = 67.103\,\mathrm{km\,s^{-1}\,Mpc^{-1}}$ from module
\texttt{h0\_diagnostic\_closure\_v28}; all other $H_0$ pathways are
classified as \emph{diagnostics}. The suite contract
(\texttt{h0\_source\_contract\_audit\_v1}, PASS) enforces this hierarchy:
only one module feeds $H_0$ into the status dashboard; the remaining pathways are
cross-checked for consistency but do not produce the headline number.
This prevents the ambiguity of a referee asking ``which $H_0$ is your claim?''

\begin{figure}[H]
\centering
\begin{tikzpicture}
\begin{groupplot}[
    group style={group size=2 by 1, horizontal sep=1.2cm},
    width=0.45\linewidth,
    height=0.35\linewidth
]
\nextgroupplot[
    ylabel={$\chi^2$},
    title={Dataset $\chi^2$ contributions (10 enabled)},
    symbolic x coords={aCov,CKM,PMNS,prim,OmDM,Omb,beta,mbb,DEgap,rUp},
    xtick=data,
    xticklabel style={rotate=35,anchor=east,font=\tiny},
    xticklabels={$\alpha$ cov, CKM, PMNS, primord., {$\Omega_\mathrm{DM}$}, {$\Omega_b$}, $\beta$, {$m_{\beta\beta}$}, DE gap, {$r_{95}$}},
    grid=major,
    ymin=0,
]
\addplot[ybar, fill=blue!50] coordinates {
    (aCov,1.8065642789580403)
    (CKM,8.50642309686793)
    (PMNS,3.4601486967025794)
    (prim,0.1092887845465857)
    (OmDM,7.56712703411243)
    (Omb,4.928781445366956)
    (beta,0.5903174784328564)
    (mbb,0.0)
    (DEgap,0.0)
    (rUp,0.0)
};

\nextgroupplot[
    ylabel={$\chi^2$},
    title={Non-zero $\chi^2$ sectors (total 27.0; 25 dof)},
    symbolic x coords={aCov,CKM,PMNS,prim,OmDM,Omb,beta},
    xtick=data,
    xticklabel style={rotate=35,anchor=east,font=\tiny},
    xticklabels={$\alpha$ cov, CKM, PMNS, primord., {$\Omega_\mathrm{DM}$}, {$\Omega_b$}, $\beta$},
    grid=major,
    ymin=0,
]
\addplot[ybar, fill=orange!60] coordinates {
    (aCov,1.8065642789580403)
    (CKM,8.50642309686793)
    (PMNS,3.4601486967025794)
    (prim,0.1092887845465857)
    (OmDM,7.56712703411243)
    (Omb,4.928781445366956)
    (beta,0.5903174784328564)
};
\end{groupplot}
\end{tikzpicture}
\caption{Likelihood breakdown. Dataset and sector $\chi^2$ contributions from the unified likelihood engine.}
\label{fig:likelihood-breakdown}
\end{figure}

\begin{figure}[H]
\centering
\begin{tikzpicture}
\begin{groupplot}[
    group style={group size=2 by 1, horizontal sep=1.2cm},
    ybar,
    width=0.48\textwidth,
    height=0.42\textwidth,
    xmin=0.5, xmax=4.5,
    xtick={1,2,3,4},
    xticklabels={$\beta_{\rm deg}$, $\lambda$, $n_s$, $A_s$},
    x tick label style={rotate=0, font=\scriptsize},
    grid=both,
]
\nextgroupplot[
    title={$\chi^2$ contributions (excluding $\alpha$)},
    ylabel={$\chi^2$},
    ymin=0, ymax=0.8,
]
\addplot coordinates {
    (1,0.5903174784328564)
    (2,0.006072422003661537)
    (3,0.021391549817205795)
    (4,0.08789723472938661)
};

\nextgroupplot[
    title={$z$-scores (excluding $\alpha$)},
    ylabel={$z$},
    ymin=-1.0, ymax=0.15,
]
\addplot coordinates {
    (1,-0.7683212078505034)
    (2,0.07792574673149778)
    (3,-0.14625850340136054)
    (4,-0.296474678057649)
};
\addplot[black!50, dashed] coordinates {(0.5,0) (4.5,0)};
\end{groupplot}
\end{tikzpicture}
\caption{Engineering-mode dashboard plots from \texttt{global\_consistency\_test}: $\chi^2$ contributions (left) and $z$-scores (right), both excluding $\alpha$.}
\end{figure}

%=============================================================================
\section{\texorpdfstring{E$_8$ Chain Certificate}{E8 Chain Certificate}}
\label{app:E8cert}
%=============================================================================

This appendix provides the complete verification data for the E$_8$ chain uniqueness 
theorem (Theorem~\ref{thm:chainunique}).

\subsection{Orbit Data Structure}

Each nilpotent orbit $\mathcal{O}$ of $\mathfrak{e}_8$ is characterized by:
\begin{itemize}[leftmargin=1.5em, topsep=2pt, itemsep=1pt]
    \item \textbf{Bala--Carter label:} Standard notation (e.g., $A_1$, $2A_1$, $A_2$, \ldots)
    \item \textbf{Dimension:} $\dim(\mathcal{O})$
    \item \textbf{Centralizer dimension:} $D(\mathcal{O}) = 248 - \dim(\mathcal{O})$
    \item \textbf{Height:} $h(\mathcal{O})$ (sum of weighted Dynkin diagram coefficients)
\end{itemize}

\subsection{The Unique Chain (Excerpt)}

The complete chain from $D = 58$ to $D = 8$:

\begin{center}
\small
\begin{tabular}{@{}clccc@{}}
\toprule
$n$ & \textbf{Orbit Label} & $D_n$ & $\dim(\mathcal{O})$ & $h$\\
\midrule
1 & $A_1$ & 58 & 190 & 2\\
2 & $2A_1$ & 56 & 192 & 4\\
3 & $3A_1$ & 54 & 194 & 6\\
$\vdots$ & $\vdots$ & $\vdots$ & $\vdots$ & $\vdots$\\
12 & (EW anchor) & 36 & 212 & 24\\
15 & (Hadronic) & 30 & 218 & 30\\
$\vdots$ & $\vdots$ & $\vdots$ & $\vdots$ & $\vdots$\\
26 & $E_8(a_1)$ & 8 & 240 & 58\\
\bottomrule
\end{tabular}
\end{center}

\textbf{Full data:} Available in machine-readable format at the replication repository 
(see Section~\ref{app:repro}).

\subsection{Verification Algorithm}

The uniqueness is verified by:
\begin{enumerate}[leftmargin=1.5em, topsep=2pt, itemsep=1pt]
    \item Constructing the complete Hasse diagram of $\mathcal{N}(\mathfrak{e}_8)$ 
        from standard tables.
    \item Computing all cover relations with $\Delta D = 2$.
    \item Running dynamic programming to find \emph{all} paths from $D = 58$ to $D = 8$.
    \item Applying the cost function $\mathcal{S} = \sum_n (\Delta^2 \ln D_n)^2$ 
        (second-difference smoothness).
    \item Verifying that the minimum-cost path is unique (cost gap to second-best $= 0$ 
        indicates tie-breaker was not needed for the optimal path).
\end{enumerate}

%=============================================================================
\section{Boundary Trace Derivation}
\label{app:boundary}
%=============================================================================

This appendix derives the block index formula $k_B = \frac{3}{2} I_1(B)$ from 
the APS index theorem.

\subsection{Setup}

Consider a 6D manifold $M^6$ with boundary $\partial M^6 = \Sigma^5$. For a 
Weyl fermion coupled to a $U(1)$ gauge field $A$, the APS index theorem gives:
\begin{equation}
    \text{ind}(\not{\!D}) = \int_{M^6} \text{ch}(F) \wedge \hat{A}(M^6) 
    - \frac{1}{2}\eta(\not{\!D}_\Sigma)
\end{equation}
where $\eta$ is the $\eta$-invariant of the boundary Dirac operator.

\subsection{Flat Connection on Boundary Cycles}

On the orientable double cover $\widetilde{M}$, the boundary decomposes as:
\begin{equation}
    \partial\widetilde{M} = C_1 \cup C_2 \cup C_T
\end{equation}
(two boundary cycles plus the twist seam).

For a \emph{flat} $U(1)$ connection on each $C_i$ with holonomy $\theta_i = \oint_{C_i} A$, 
the $\eta$-invariant contribution becomes:
\begin{equation}
    \eta_i \propto \sum_{\text{fermions}} q^2 \cdot \theta_i
\end{equation}
where $q$ is the $U(1)$ charge.

\subsection{\texorpdfstring{Spectral Trace Coefficient and $\bone$}{Spectral Trace Coefficient and b1}}
The regulated determinant of the hypercharge-coupled Dirac operator contains a logarithmic
term which we normalize as:
\begin{equation}
    \ln\det(\not{D}_Y) \supset \frac{1}{2\pi}\,b_Y\,\ln\frac{\Lambda}{\mu}\;\int d^4x\sqrt{-g}\,F_{\mu\nu}F^{\mu\nu},
    \qquad
    b_Y=\frac{2}{3}\sum_{\text{Weyl}} d_3 d_2\,Y^2+\frac{1}{3}\sum_{\text{scalars}} d_3 d_2\,Y^2.
\end{equation}
We define the spectral index $b_Y$ as the quadratic charge trace entering this term in our
normalization; the overall prefactor is part of the declared coefficient policy (Appendix~C)
and differs by a fixed normalization from the conventional $1/(16\pi^2)$ Seeley--DeWitt prefactor.
With $Q=T_3+Y$ and the SM field content, the explicit arithmetic is:
\begin{center}
\begin{tabular}{@{}lcccccc@{}}
\toprule
Field & $Q$ & $L$ & $u_c$ & $d_c$ & $e_c$ & $H$ \\
\midrule
Contribution to $b_Y$ & $\tfrac{1}{3}$ & $1$ & $\tfrac{8}{3}$ & $\tfrac{2}{3}$ & $2$ & $\tfrac{1}{6}$ \\
\bottomrule
\end{tabular}
\end{center}
so $b_Y=41/6$ and $\bone=\frac{3}{5}b_Y=\frac{41}{10}$.
Thus $\bone$ appears as a \emph{spectral index coefficient},
independent of IR thresholds, matching its role in the CFE.
\textbf{Policy note:} the overall prefactor is fixed by the declared coefficient scheme
(background-field gauge, $\overline{\mathrm{MS}}$) and is treated as part of the CFE coefficient policy
until a full background-field derivation is provided.

\subsection{\texorpdfstring{Derivation of $k_B$}{Derivation of kB}}

\textbf{Step 1:} Each boundary cycle contributes a factor proportional to $I_1 = \sum q^2$.

\textbf{Step 2:} Three cycles contribute additively: factor 3.

\textbf{Step 3:} The double cover counts each physical state once across two sheets. 
The effective weight is halved: factor $1/2$.

\textbf{Result:}
\begin{equation}
    k_B = 3 \times \frac{1}{2} \times I_1(B) = \frac{3}{2} I_1(B)
\end{equation}

\textbf{Consistency check:} For the electroweak block, 
$I_1^{\text{EW}} = \frac{41}{48}$ (SM hypercharge sum), giving 
$k_{\text{EW}} = \frac{3}{2} \times \frac{41}{48} = \frac{41}{32}$, 
as stated in Table~\ref{tab:blocks}.

%=============================================================================
